\section{Introduction}
\label{sec:intro}

Traffic engineering (TE) decides which paths carry which traffic so that a
backbone delivers its demand without congestion, SLA violations or
needless reconfiguration. In MPLS networks the decision is concrete: each
traffic trunk rides one label-switched path, and an operator or controller
moves trunks between a few precomputed candidates as load
changes~\citep{rfc2702,rfc3209}. Because traffic varies over the day and
failures arrive unannounced, reinforcement learning (RL) is an appealing
tool, and deep RL has been applied to routing and TE in many
forms~\citep{valadarsky2017learning,xu2018drlte,zhang2020cfrrl,almasan2022drlgnn,bernardez2023magnneto}.

RL earns its machinery---value functions, discounting, credit assignment
over long horizons---only if today's decision changes tomorrow's options in
ways that a myopic decision rule would get wrong. Whether that is true of TE
is rarely measured. The strongest recent TE systems argue the opposite for
the standard setting, in which routing is re-optimized from scratch every
interval and therefore has no memory: Teal reduces its RL objective to a
one-step return~\citep{xu2023teal}, and DOTE reports that direct optimization
beats an RL-based scheme by a wide margin~\citep{perry2023dote}. Operational
MPLS-TE is different in exactly the property those arguments rely on. LSP
moves are incremental and signalled, controllers limit how many trunks change
at once, stability timers prevent a trunk from moving again too soon, and
protected traffic may move only where capacity is guaranteed. The current
assignment is then part of the state, a move constrains the future, and the
decision problem is a genuine Markov decision process.

This paper asks, for such a formulation, \emph{whether sequential RL is
needed, and why or why not}. We study a flow-level MPLS backbone (18 routers,
64 directed links, 17 demands with SLAs, four candidate LSPs each) in which a
controller may move one demand per five-minute interval, moves cost
reconfiguration effort, a three-interval hold-down follows each move, and an
action mask enforces failed links, hold-downs and protected-class capacity.
We compare MaskablePPO~\citep{schulman2017ppo,huang2022masking} with a masked
neural contextual bandit that regresses only the immediate reward of each
action, under identical observations, masks, budgets and paired traffic.

The comparison was first run as a governed study by the project that
developed the environment, which found the bandit ahead. We treat that study
as a claim to be tested. Our contributions are:

\begin{itemize}
\item \textbf{A reproduction and its lesson.} Retraining all of the closed
study's runs on different hardware with byte-identical training traffic, the
bandit's result reproduces and PPO's does not: floating-point differences
alone move PPO's per-root holdout return by up to \ReproMaxShiftPPO{} points
and the bandit's by at most \ReproMaxShiftBandit{} (\S\ref{sec:repro}).
\item \textbf{A controlled comparison with a non-learning optimizer.} On
\TestBanditRoots{} training roots and 140 disjoint test episodes the bandit
beats PPO on every root and is statistically indistinguishable from a
per-interval min-max-utilization MILP controller that does not learn at all
(\S\ref{sec:rq1}).
\item \textbf{A mechanism.} PPO's deficit is dominated by reconfiguration
cost: on most roots it settles into flipping a demand at every hold-down
expiry, a small immediate cost that its long-horizon advantage estimates fail
to attribute (\S\ref{sec:analysis}). Horizon and algorithm are separated by
varying the discount within each learner family (\S\ref{sec:rq2}).
\item \textbf{A measured sequentiality diagnostic.} Clairvoyant rollouts show
that the problem is far from myopic per decision under the true future, but
that almost all of the non-myopic value is anticipation of exogenous traffic
change; with traffic held fixed, the myopic choice is the 24-interval choice
in \SeqFrozenAgreeHTwentyFour\,\% of states (\SeqLiveAgreeHTwentyFour\,\% under the true future) (\S\ref{sec:sequentiality}).
\item \textbf{A regime where sequential RL is necessary.} Delaying the effect
of a move by one interval---a controlled form of actuation latency---removes
the bandit's signal by construction; the bandit collapses while PPO is
unaffected, reversing the ranking (\S\ref{sec:rq5}).
\end{itemize}

We make no algorithmic or theoretical claim. All code, configurations,
per-episode results and the scripts that produce every number, table and
figure in this paper are part of the artifact.
